\chapter{Targets, Labels and Sample Weights}\label{ch:ml:targets-labels-and-sample-weights}

A classifier is trained to say whether an asset will be up ten days from now, and it does so a little better than a
coin. The trader who will use it asks what it does about the stop-loss: every position is closed if it loses one
ten-day volatility, and taken profit if it gains one. On twenty synthetic assets over ten years, one trade in five hits
its stop before day ten; the label the classifier learned never saw the stop. And the \omterm{def:ml:why-financial-machine-learning-is-different:sets}{training set} that looked like
49\,000 examples holds, once the overlap of ten-day windows is counted, the information of 4\,900. This chapter builds
labels that describe the trade, measures how much labels overlap, and turns the overlap into \omterm{def:ml:targets-labels-and-sample-weights:weights}{sample weights} and
honest statistics.

\section{Fixed-horizon labels and their flaws}

The data of the chapter come from \texttt{firm.mlsynth.series}: twenty assets, 2\,520 days, daily volatility 1.5\% with
GARCH clustering and Student-$t$ shocks, and a planted drift, persistent with a half-life of 20 days, whose correlation
with the next day's return is 0.04. The features, known at each close, are a noisy reading of the drift, the 5- and
20-day past returns, a ratio of short to long realised volatility and three pure-noise columns. An event is every day
of every asset; the primary trade is long when the drift reading is positive and short otherwise.

\begin{definition}[Fixed-horizon label]\label{def:ml:targets-labels-and-sample-weights:fixed}
Given an event at bar $t_0$ and a horizon $h$, the \emph{fixed-horizon label}\index{fixed-horizon label} is
$\sign(\sum_{s=t_0+1}^{t_0+h}R_s)$, or the return itself for regression, settled at $t_1 = t_0 + h$.
\end{definition}

The \omterm{def:ml:targets-labels-and-sample-weights:fixed}{fixed-horizon label} is Book~7's forward return (chapter~6) turned into a class. It has three flaws. It ignores the
path: a position that would have been stopped out on day three and a position that ended up on day ten get the same
label if the price came back. It ignores the volatility: a 1\% move is noise on a volatile day and news on a quiet one.
And it ends at a fixed time, so the labels of consecutive events overlap in $h - 1$ of their $h$ bars. The first two
are fixed by the barrier label below; the third by counting (\cref{sec:ml:labels:overlap}).

\section{Barrier labels and meta-labelling}

\begin{definition}[Triple-barrier label]\label{def:ml:targets-labels-and-sample-weights:triple}
For an event at $t_0$ with position side $\pm1$, widths $w^+, w^- > 0$ and a horizon $h$, the \emph{triple-barrier
label}\index{triple-barrier label} settles at the first bar $t_1$ at which the position's cumulative log return reaches
$+w^+$ (profit-taking), $-w^-$ (stop-loss), or $t_1 = t_0 + h$ (the vertical barrier); it is the sign of the position's
return at $t_1$. The widths are usually scaled by the volatility known at $t_0$: $w = k\hat\sigma_{t_0}\sqrt h$
(López de Prado, 2018).
\end{definition}

With $k = 1$ and $h = 10$, the long-or-short trade of the chapter hits its profit barrier first 22.7\% of the time, its
stop 20.3\% of the time, and the vertical barrier otherwise; a label lasts 8.3 days on average instead of 10.
\Cref{fig:ml:labels:path} shows one event of the first asset: a long position stopped on day three at $-5.04\%$ (the
width was 4.87\%), when the price went on to end day ten 1.84\% higher. The \omterm{def:ml:targets-labels-and-sample-weights:fixed}{fixed-horizon label} calls that trade a
winner. It is one case of many: of the trades stopped out, 5.3\% end positive at day ten. Signs disagree for 1.9\% of
all events; what differs everywhere is the size, since the barrier label's returns are truncated at the barriers as the
trade's are.

\begin{omfigure}
\begin{tikzpicture}
\begin{axis}[width=10.5cm,height=5.6cm,xlabel={days after the event},ylabel={cumulative return (\%)},
  tick label style={font=\small},label style={font=\small},xmin=-0.3,xmax=10.3,ymin=-7,ymax=6.5,
  xtick={0,1,2,3,4,5,6,7,8,9,10},grid=major,grid style={gray!25},table/col sep=comma]
\addplot[omDef,very thick,mark=*,mark size=1.6pt] table[x=day,y=ret]{figdata/ml/02-targets-labels-and-sample-weights/path.csv};
\addplot[omProp,thick,dashed] table[x=day,y=up]{figdata/ml/02-targets-labels-and-sample-weights/path.csv};
\addplot[omThm,thick,dashed] table[x=day,y=down]{figdata/ml/02-targets-labels-and-sample-weights/path.csv};
\draw[black!70,thick] (axis cs:10,-7) -- (axis cs:10,6.5);
\node[font=\footnotesize,omProp,anchor=south west] at (axis cs:0.2,4.87) {profit-taking};
\node[font=\footnotesize,omThm,anchor=north] at (axis cs:8.2,-5.0) {stop-loss};
\node[font=\footnotesize,black!70,anchor=east] at (axis cs:9.9,5.6) {vertical barrier};
\draw[omThm,thick] (axis cs:3,-5.04) circle[radius=4pt];
\node[font=\footnotesize,omThm,anchor=west] at (axis cs:3.3,-6.2) {stopped on day 3};
\end{axis}
\end{tikzpicture}

\omcaption{One event of asset~0, a long position: the barriers are one ten-day volatility, 4.87\%, known at the event.
The stop is touched on day three; the ten-day return ends at $+1.84\%$. Data: \texttt{firm.mlsynth.series}, seed~1,
through \texttt{firm.labeling.triple\_barrier}.}\label{fig:ml:labels:path}
\end{omfigure}

The side of a trade and its size are different questions. A primary model (a rule, a researcher's signal, a
discretionary call) may be good at the side and poor at knowing when it is right.

\begin{definition}[Meta-labelling]\label{def:ml:targets-labels-and-sample-weights:meta}
\emph{Meta-labelling}\index{meta-labelling} labels each event of a primary model with 1 if the trade in the primary
model's direction made money (for example, its triple-barrier return is positive) and 0 otherwise, and trains a
secondary model on these labels to estimate the probability that the primary call is right; the probability is then
used to filter or size the trade (López de Prado, 2018; Joubert, 2022).
\end{definition}

\omterm{def:ml:targets-labels-and-sample-weights:meta}{Meta-labelling} the chapter's trade with gradient-boosted trees (\cref{fig:ml:labels:meta}), trained on the first 60\%
of the dates and scored on the last 40\%, raises the hit rate of the trades it keeps from 53.0\% to 54.5\% and the
average return per trade from 0.064 to 0.097 barrier widths when it keeps the third of the trades it is most confident
about. The overlap-adjusted $t$-statistic falls from 3.48 to 3.19: the kept trades are better, and there are fewer of
them. \omterm{def:ml:targets-labels-and-sample-weights:meta}{Meta-labelling} reallocates risk to the better trades; it does not create a signal the primary model lacks.

\begin{omfigure}
\begin{tikzpicture}
\begin{groupplot}[group style={group size=2 by 1,horizontal sep=1.7cm},width=6.6cm,height=5.2cm,
  xlabel={\small trades kept (\%)},tick label style={font=\footnotesize},label style={font=\small},x dir=reverse,
  xmin=5,xmax=105,grid=major,grid style={gray!25},table/col sep=comma]
\nextgroupplot[ylabel={mean per trade (widths)},ymin=0.05,ymax=0.13,
  yticklabel style={/pgf/number format/fixed,/pgf/number format/precision=2},scaled y ticks=false]
\addplot[omDef,very thick,mark=*] table[x=share,y=mean]{figdata/ml/02-targets-labels-and-sample-weights/meta.csv};
\nextgroupplot[ylabel={overlap-adjusted $t$},ymin=2,ymax=3.8]
\addplot[omThm,very thick,mark=square*] table[x=share,y=t]{figdata/ml/02-targets-labels-and-sample-weights/meta.csv};
\end{groupplot}
\end{tikzpicture}

\omcaption{\omterm{def:ml:targets-labels-and-sample-weights:meta}{Meta-labelling} the primary trade: keeping the trades whose estimated probability of success exceeds 0.50 to
0.55 (100\% = the primary model alone). Left, the average return per trade in barrier widths; right, the $t$-statistic
with the number of trades adjusted for overlap. Out of sample, last 40\% of the dates. Data: \texttt{ml\_labels.meta}.}\label{fig:ml:labels:meta}
\end{omfigure}

\section{Overlap, concurrency and uniqueness}\label{sec:ml:labels:overlap}

\begin{definition}[Label concurrency, average uniqueness]\label{def:ml:targets-labels-and-sample-weights:uniqueness}
Label $i$ spans the bars $(t_{0,i}, t_{1,i}]$. The \emph{label concurrency}\index{label concurrency} $c_t$ of bar $t$
is the number of labels whose span contains it. The \emph{average uniqueness}\index{average uniqueness} of label $i$ is
$\bar u_i = \frac{1}{t_{1,i} - t_{0,i}}\sum_{t\in(t_{0,i}, t_{1,i}]}1/c_t$: the share of its information it holds alone.
\end{definition}

\Cref{fig:ml:labels:toy} is the smallest example: three labels over seven bars, two of them overlapping in two bars.
Book~7 named the phenomenon (label overlap, chapter~20) and removed its leak from cross-validation by purging; here it is
measured, because it changes how much the sample is worth.

\begin{omfigure}
\begin{tikzpicture}[x=1.15cm,y=0.62cm,font=\small]
\foreach \b in {1,2,3,4,5,6,7} { \draw[gray!50] (\b-0.5,-0.2) -- (\b-0.5,4.4); \node at (\b,4.1) {\b}; }
\draw[gray!50] (7.5,-0.2) -- (7.5,4.4);
\node[anchor=east] at (0.3,4.1) {bar};
\draw[omDef,line width=4pt] (0.55,3) -- (3.45,3); \node[anchor=east] at (0.3,3) {label A};
\draw[omProp,line width=4pt] (1.55,2) -- (4.45,2); \node[anchor=east] at (0.3,2) {label B};
\draw[omThm,line width=4pt] (5.55,1) -- (7.45,1); \node[anchor=east] at (0.3,1) {label C};
\node[anchor=east] at (0.3,0) {$c_t$};
\foreach \b/\c in {1/1,2/2,3/2,4/1,5/0,6/1,7/1} { \node at (\b,0) {\c}; }
\node[anchor=west,align=left,font=\footnotesize] at (7.8,2) {$\bar u_A = (1+\tfrac12+\tfrac12)/3 = 0.67$\\
$\bar u_B = (\tfrac12+\tfrac12+1)/3 = 0.67$\\ $\bar u_C = 1$};
\end{tikzpicture}

\omcaption{Concurrency and \omterm{def:ml:targets-labels-and-sample-weights:uniqueness}{average uniqueness} on seven bars: labels A and B share bars 2 and 3; C is alone. The three
labels are worth $0.67 + 0.67 + 1 = 2.33$ independent labels, the number of bars they cover divided by their length only
when all lengths are equal (\cref{prop:ml:labels:sum}).}\label{fig:ml:labels:toy}
\end{omfigure}

\begin{proposition}[How much overlapping labels are worth]\label{prop:ml:labels:sum}
\begin{enumerate}
\item If every label spans $h$ bars, $\sum_i\bar u_i = N_{\mathrm{cov}}/h$, where $N_{\mathrm{cov}}$ is the number of
bars covered by at least one label. With an event on every bar, $\bar u_i = 1/h$ away from the ends.
\item If returns are uncorrelated with variance $\sigma^2$ and the $n$ labels are the overlapping $h$-bar sums of one
series, the sample mean of the labels has variance about $h^2\sigma^2/n$, $h$ times the $h\sigma^2/n$ that treating them
as independent assumes: a $t$-statistic computed on $n$ rows is inflated by about $\sqrt h$.
\end{enumerate}
\end{proposition}

\begin{proof}
(1) $\sum_i\bar u_i = \frac1h\sum_i\sum_{t\in i}1/c_t = \frac1h\sum_t\sum_{i\ni t}1/c_t = \frac1h\#\{t: c_t\ge1\}$,
since the inner sum has $c_t$ terms equal to $1/c_t$. (2) Away from the ends each return enters $h$ of the $n$ sums, so
the mean of the sums is about $\frac hn\sum_tR_t$, of variance $\frac{h^2}{n^2}\,n\sigma^2 = h^2\sigma^2/n$.
\end{proof}

On the chapter's data the \omterm{def:ml:targets-labels-and-sample-weights:fixed}{fixed-horizon labels} have a concurrency of exactly 10 and an \omterm{def:ml:targets-labels-and-sample-weights:uniqueness}{average uniqueness} of 0.100:
48\,980 labels are worth 4\,916, as part 1 predicts ($20\times2\,458/10$). The barrier labels, shorter, have a uniqueness
of 0.123 and are worth 6\,040. The primary trade's average return has a $t$-statistic of 9.4 computed on the rows and
3.3 with the number of trades replaced by the sum of uniqueness: the ratio, 2.85, is $\sqrt{1/0.123}$.

\section{Sample weights and the sequential bootstrap}

\begin{definition}[Sample weight, time-decay weight]\label{def:ml:targets-labels-and-sample-weights:weights}
A \emph{sample weight}\index{sample weight} $\omega_i\ge0$ scales observation $i$'s term in the empirical loss,
$\mathcal L_n(\theta) = \sum_i\omega_i\ell(y_i, f_\theta(x_i))/\sum_i\omega_i$. Three are common on market data:
uniqueness weights $\omega_i = \bar u_i$; return-attribution weights, $\omega_i\propto|\sum_{t\in i}R_t/c_t|$, which
favour labels whose moves they own; and a \emph{time-decay weight}\index{time-decay weight}, which falls linearly from 1
for the newest label to a floor for the oldest, in units of cumulative uniqueness, so that old data count less.
\end{definition}

\begin{definition}[Sequential bootstrap]\label{def:ml:targets-labels-and-sample-weights:seqboot}
The \emph{sequential bootstrap}\index{sequential bootstrap} draws labels one at a time, with replacement, each with
probability proportional to the \omterm{def:ml:targets-labels-and-sample-weights:uniqueness}{average uniqueness} it would have given the labels already drawn; overlapping labels
become less likely once one of them is in the sample.
\end{definition}

The case for these tools is that a bagged ensemble whose bootstrap samples are full of near-duplicates grows correlated
trees, each fitted to the same few independent observations. Drawing 40 of the first 400 ten-day labels of one asset
(about their effective number), the plain bootstrap's samples have an \omterm{def:ml:targets-labels-and-sample-weights:uniqueness}{average uniqueness} of 0.65 and the \omterm{def:ml:targets-labels-and-sample-weights:seqboot}{sequential
bootstrap}'s 0.69. Whether it matters for prediction is an empirical question, and on this data the answer is no
(\cref{tab:ml:labels:bagging}): bagged trees trained on plain bootstrap samples, on samples as small as the \omterm{def:ml:targets-labels-and-sample-weights:uniqueness}{average
uniqueness} times the sample size, and on those samples weighted by uniqueness have the same out-of-sample AUC in three
simulated markets, within $\pm0.007$ of each other and in no consistent order. The trees were regularised (at least 50
labels per leaf, 5 in the small samples), and none of the features identifies a date, which is what lets
near-duplicates leak in the first place. Where overlap does bite without fail is in every statistic that counts rows:
standard errors, $t$-statistics, bootstrap intervals and out-of-bag scores.

\begin{table}[htbp]
\centering\small
\begin{tabular}{@{}lcccc@{}}
\toprule
bootstrap samples & market 1 & market 2 & market 3 & mean\\
\midrule
plain, full size & 0.504 & 0.517 & 0.509 & 0.510\\
average-uniqueness size & 0.511 & 0.512 & 0.506 & 0.510\\
same, weighted by uniqueness & 0.506 & 0.511 & 0.510 & 0.509\\
\bottomrule
\end{tabular}
\caption{Out-of-sample AUC of 30 bagged trees predicting the sign of the ten-day return, trained on the first 60\% of
the dates (purged) of three simulated markets. Data: \texttt{ml\_labels.bagging\_table}.}\label{tab:ml:labels:bagging}
\end{table}

\begin{method}[Labels and weights for a new study]\label{met:ml:labels:method}
\begin{enumerate}
\item Write the trade first (side, entry, exits, horizon), then the label that settles it: barriers at the trade's
stop and target, scaled by volatility known at the event.
\item Store every label with its $t_0$ and $t_1$; compute concurrency and uniqueness per asset.
\item Report the effective sample $\sum_i\bar u_i$ beside the number of rows, and compute every standard error with it
(or with a block bootstrap of Book~4, chapter~13).
\item Try uniqueness and \omterm{def:ml:targets-labels-and-sample-weights:weights}{time-decay weights} as \omterm{def:ml:why-financial-machine-learning-is-different:sets}{hyperparameters}, on validation folds that are purged (chapter~3); keep
them only if they help there.
\end{enumerate}
\end{method}

\section{Tutorial: labels that describe the trade}\label{tut:ml:targets-labels-and-sample-weights:tut}

\begin{tutorial}
\textbf{Goal.} Label twenty assets' daily events with fixed-horizon and \omterm{def:ml:targets-labels-and-sample-weights:triple}{triple-barrier labels}, measure overlap, compare
bootstrap schemes and meta-label the primary trade. \textbf{End state:} the numbers of the chapter,
\cref{fig:ml:labels:path,fig:ml:labels:meta}, \cref{tab:ml:labels:bagging}.
\begin{tutsteps}
\item \textbf{Barrier labels}: the first touch of either barrier, or the vertical one.
\omcode{code/firm/labeling/firm_labeling.py}{41}{70}{The triple-barrier label.}\label{lst:ml:labels:triple}
\item \textbf{Concurrency and uniqueness} by cumulative sums.
\omcode{code/firm/labeling/firm_labeling.py}{77}{92}{Concurrency and average uniqueness.}\label{lst:ml:labels:uniq}
\item \textbf{The \omterm{def:ml:targets-labels-and-sample-weights:seqboot}{sequential bootstrap}}: each draw favours the labels that overlap least with those already drawn.
\omcode{code/firm/labeling/firm_labeling.py}{110}{123}{The sequential bootstrap.}\label{lst:ml:labels:seq}
\item \textbf{Run} \texttt{ml\_labels.label\_stats()}, \texttt{bagging\_table()}, \texttt{bootstrap\_uniqueness()},
\texttt{meta()} and \texttt{fig\_labels.py}.
\end{tutsteps}
\textbf{What to change next.} Set the barriers at two volatilities and see how the share of vertical-barrier exits and
the uniqueness change; meta-label with a logistic regression instead of boosted trees.
\end{tutorial}

\section{Build: labels and sample weights}\label{bld:ml:targets-labels-and-sample-weights:labeling}

\begin{build}
\bfield{Purpose} Every study labels its events the way the desk trades them, and knows how many independent
observations it holds.
\bfield{Interface} \texttt{fixed\_horizon(r, t0, h)}, \texttt{triple\_barrier(r, t0, up, down, h, side)},
\texttt{vol\_widths(sigma, t0, h, k)}, \texttt{meta\_labels(side, ret)}, \texttt{concurrency(t0, t1, n)},
\texttt{average\_uniqueness(t0, t1, n)}, \texttt{attribution\_weights(t0, t1, r, n)}, \texttt{time\_decay(u, last)},
\texttt{sequential\_bootstrap(t0, t1, n, size, rng)}.
\bfield{Rules} Every label carries $t_0$ and $t_1$; bar $t_0$'s close is the decision time and the label uses bars
$t_0 + 1$ to $t_1$; barrier widths use only information known at $t_0$.
\bfield{Acceptance tests} \texttt{code/firm/labeling/tests/}: labels and touches by hand, for a long and a short;
concurrency, uniqueness, attribution and decay weights on a three-label example; the \omterm{def:ml:targets-labels-and-sample-weights:seqboot}{sequential bootstrap} raises the
\omterm{def:ml:targets-labels-and-sample-weights:uniqueness}{average uniqueness} of small samples.
\bfield{Stretch} Trend-scanning labels (the horizon with the most significant trend); labels for market-making fills
(chapter~29); a vectorised barrier search over many events at once.
\end{build}

\begin{omsources}
\item M.~López de Prado, \emph{Advances in Financial Machine Learning}, Wiley, 2018, chapters 3--4.
\item J.~F.~Joubert, ``Meta-labeling: theory and framework'', \emph{Journal of Financial Data Science} 4(3), 2022.
\item L.~Breiman, ``Bagging predictors'', \emph{Machine Learning} 24, 1996.
\end{omsources}

\section{Exercises}

\begin{exercise}[$\star$]\label{exo:ml:targets-labels-and-sample-weights:1}
Labels of 5 bars start on every bar of a 1\,000-bar series (the last ones truncated at the end). About how many
independent labels do they amount to, and what is their \omterm{def:ml:targets-labels-and-sample-weights:uniqueness}{average uniqueness} away from the ends?
\end{exercise}

\begin{exercise}[$\star$]\label{exo:ml:targets-labels-and-sample-weights:2}
Four labels span bars $(0,2]$, $(1,4]$, $(2,4]$ and $(6,8]$. Write the concurrency of bars 1 to 8 and each label's
\omterm{def:ml:targets-labels-and-sample-weights:uniqueness}{average uniqueness}.
\end{exercise}

\begin{exercise}[$\star$]\label{exo:ml:targets-labels-and-sample-weights:3}
The daily volatility known at an event is 1.2\%. What are the barrier widths for $k = 1.5$ and a horizon of 16 days?
\end{exercise}

\begin{exercise}[$\star\star$]\label{exo:ml:targets-labels-and-sample-weights:4}
A strategy's average trade has a $t$-statistic of 6.1 computed on 25\,000 overlapping 20-day trades opened every day.
What $t$-statistic should be reported, and why?
\end{exercise}

\begin{exercise}[$\star\star$]\label{exo:ml:targets-labels-and-sample-weights:5}
Why does \omterm{def:ml:targets-labels-and-sample-weights:meta}{meta-labelling} raise the average return per trade but not the $t$-statistic of the whole book, in
\cref{fig:ml:labels:meta}? When would it raise both?
\end{exercise}

\begin{exercise}[$\star\star$]\label{exo:ml:targets-labels-and-sample-weights:6}
\emph{Find the flaw.} ``Our barrier widths are set at one standard deviation of each asset's daily returns over the whole
sample, times the square root of the horizon.''
\end{exercise}

\begin{exercise}[$\star\star\star$]\label{exo:ml:targets-labels-and-sample-weights:7}
\emph{Coding.} Rebuild \texttt{ml\_labels.events} with barriers at two ten-day volatilities ($k = 2$). Report the shares
of profit, stop and vertical exits and the \omterm{def:ml:targets-labels-and-sample-weights:uniqueness}{average uniqueness} of the barrier labels, and explain the change.
\end{exercise}

\begin{exercise}[$\star\star\star$]\label{exo:ml:targets-labels-and-sample-weights:8}
With \omterm{def:ml:targets-labels-and-sample-weights:weights}{time-decay weights} falling linearly from 1 to 0.5 in cumulative uniqueness, and every label of uniqueness 0.1, what
is the weight of the label halfway through the sample? Show that the weights' mean is 0.75 whatever the uniqueness if it
is constant, and say what changes when it is not.
\end{exercise}

\section{Problem: Five Hundred Overlapping Days}

\begin{problem}[{Weekend problem --- what a training set is worth}]\label{pb:ml:targets-labels-and-sample-weights:1}
The chapter's twenty assets, their daily events and the primary long-or-short trade.

\medskip\noindent\textbf{Part I --- Two labels.}
\begin{enumerate}
\item How often does the trade hit its profit barrier, its stop, the vertical barrier?
\item How long does a barrier label last on average, and why is it shorter than ten days?
\item In \cref{fig:ml:labels:path}, what do the fixed-horizon and the barrier label say, and which describes the trade?
\item What share of all events have labels of opposite signs, and why is that not the whole difference?
\end{enumerate}

\medskip\noindent\textbf{Part II --- Overlap.}
\begin{enumerate}[resume]
\item What are the concurrency and the \omterm{def:ml:targets-labels-and-sample-weights:uniqueness}{average uniqueness} of the \omterm{def:ml:targets-labels-and-sample-weights:fixed}{fixed-horizon labels}?
\item How many independent labels are the 48\,980 worth, by \cref{prop:ml:labels:sum}? And the barrier labels?
\item Give the trade's $t$-statistic on rows and adjusted for overlap, and explain their ratio.
\item Which statistics of a study are wrong if overlap is ignored?
\end{enumerate}

\medskip\noindent\textbf{Part III --- Weights and bootstraps.}
\begin{enumerate}[resume]
\item What \omterm{def:ml:targets-labels-and-sample-weights:uniqueness}{average uniqueness} do plain and \omterm{def:ml:targets-labels-and-sample-weights:seqboot}{sequential bootstrap} samples of 40 labels have?
\item What AUC do the three bagging schemes reach, market by market?
\item Why did uniqueness make no difference to the AUC here?
\item Design a setting where it would.
\end{enumerate}

\medskip\noindent\textbf{Part IV --- \omterm{def:ml:targets-labels-and-sample-weights:meta}{Meta-labelling} and the verdict.}
\begin{enumerate}[resume]
\item What do the kept trades gain at a threshold of 0.52, and what does the book lose?
\item Why is \omterm{def:ml:targets-labels-and-sample-weights:meta}{meta-labelling} sizing, not alpha?
\item Would you use the probability to size rather than to filter? How?
\item State the \emph{named result}: the effective number of observations against the nominal one, the $t$-statistic
before and after, and the AUC gain of uniqueness weighting over plain bagging with its spread.
\item What would the effective sample be with 20-day labels?
\item How should a research log record the size of a \omterm{def:ml:why-financial-machine-learning-is-different:sets}{training set}?
\item What does purging (Book~7) have to do with uniqueness?
\item In one sentence: what is a label?
\end{enumerate}
\end{problem}

\section{Interview questions}

\begin{interviewq}[$\star$ \iqroles{researcher, mle}]\label{iq:ml:targets-labels-and-sample-weights:1}
What is wrong with labelling each day by the sign of the next 20 days' return?
\end{interviewq}

\begin{interviewq}[$\star$ \iqroles{researcher}]\label{iq:ml:targets-labels-and-sample-weights:2}
Explain the triple-barrier method in two minutes.
\end{interviewq}

\begin{interviewq}[$\star\star$ \iqroles{researcher, mle}]\label{iq:ml:targets-labels-and-sample-weights:3}
You have a million rows of 10-day labels sampled daily from 400 stocks over ten years. How many independent observations
do you have, roughly, and how does that change your model choice?
\end{interviewq}

\begin{interviewq}[$\star\star$ \iqroles{researcher, trader}]\label{iq:ml:targets-labels-and-sample-weights:4}
What is \omterm{def:ml:targets-labels-and-sample-weights:meta}{meta-labelling}, and when is it useful?
\end{interviewq}

\begin{interviewq}[$\star\star$ \iqroles{mle}]\label{iq:ml:targets-labels-and-sample-weights:5}
A random forest's out-of-bag score on overlapping labels is much better than its score on a later test period. Why?
\end{interviewq}

\begin{interviewq}[$\star\star\star$ \iqroles{researcher}]\label{iq:ml:targets-labels-and-sample-weights:6}
Show that the mean of $n$ overlapping $h$-period sums of white noise has about $h$ times the variance you would compute
if you treated them as independent.
\end{interviewq}
